\definepapersize[seferpage][width=170mm,height=240mm]
\setuppapersize[seferpage]

\setuplayout[
  topspace=11mm,
  bottomspace=9mm,
  backspace=14mm,
  width=142mm,
  height=220mm,
]

\mainlanguage[he]
\setupalign[r2l,hz,hanging]
\definefontfamily[hebrewfont][rm][David CLM]
\setupbodyfont[hebrewfont,12pt]

\starttext

{\bf\tfx ח \hfill {\tfd שפע} \quad שער גמילות חסדים \hfill שלמה}

\blank[big]

{\bf מעלה את הקב"ה, הוא על הטובות שמשפיע לבני ישראל, כי כשרואים בטובת ישראל מספרים גודל שבחו, לפי שבזה ניכר תוקף חסדו וכח גבורתו.}

\blank[small]
\midaligned{\tfx מהות ה'חסיד' - מתפלל ומיטיב לזולתו}
\blank[small]

{\bf ו. לעומת זאת, המרחם על חברו, משתתף בצערו ומתפלל בעדו, הרי הוא נחשב חלקו ונחלתו של הקב"ה"א. ולא ייקרא האדם בתואר 'חסיד', שנקרא כן על שם החסד שגומל, אלא אם כן הוא דורש טובת זולתו ומבקש עליו רחמים; שכן מי שמבקש, מתפלל ומשפיע רק על עצמו, אין מעשיו נחשבים לחסד כלל: [ג].}

\blank[medium]
\midaligned{◆ ◆ ◆ ◆ ◆ ◆ ◆ ◆ ◆ ◆ ◆ ◆ ◆ ◆ ◆ ◆ ◆ ◆ ◆ ◆}
\blank[medium]

\midaligned{{\bf ◆ מקור השפע ◆} \hfill {\bf ◆ צינור השפע ◆}}
\blank[small]

\startcolumns[n=2,distance=4mm,rule=on]
'ישמח ה' במעשיו' (תהלים קד לא), ונגד כבודו יתברך שמו לא היה שוה בעיניו כלל צער של בני ישראל. כך רצון ה' יתברך אינו בזה, כי חפצו ורצונו שיהיה מנהיג הדור כדי לעורר רחמים על בני ישראל ולא לחוס על כבוד הבורא לבד. וזה הענין שמצינו ברבי ישמעאל בן אלישע (ברכות ז א) שנכנס לפני ולפנים והיה מתפלל על בני ישראל.

\column

[ג] מסופר על רבינו הרה"ק בעל התפארת שלמה מראדומסק זיע"א שבשבת פרשת נח היה מסתגר בחדרו ולא היה יוצא לערוך שלחנו הטהור, באמרו על נח 'ער איז געווענזען א צדיק אין פעלץ'.
\stopcolumns

\blank[small]

{\tfx וזה היה ג"כ חסרון של נדב ואביהוא 'בקרבתם לפני ה'', פירוש שהיו הולכים תמיד רק במדריגה זו לפני ה' - לחוס רק על כבודו לבד, למסור נפש עבורי ית"ש, לכך זימותו', כי זה הדרך הוא עצמו, אבל לטובת הדור ולהנחיגם צריך הצדיק להיות בעל הרחמים לבקש טובתם ולעורר רחמים חסדים, ומאשר לא היה להם הנהגה בעוה"ז לכן נסתלקו מעוה"ז.

בפרשת פנחס (ד"ה החתן): "יש שני מיני צדיקים, האחד, הוא המייחל ומצפה על גדולת הבורא ב"ה יהי כבוד ה' לעולם בביאת משיח במהרה בימינו... אולם יש בחינה אחרת בצדיק יסוד עולם, להיות דורש טוב לעמו, אשר בזה בחר ה', כי חפצו ורצונו יתברך שמו להיות צדיק כזה להשפיע טוב לבני ישראל".}

\stoptext
